\section{La misión como autómata de fases}
\label{sec:mision}

Los resultados anteriores describen mecanismos —reclutar, certificar, encaminar,
frenar— sin decir cómo se encadenan. Esta sección los organiza en un autómata
con guardas y muestra que las tres fases no son tres sistemas distintos: son el
mismo objeto con conjuntos admisibles anidados.

\subsection{Fases y guardas}

\begin{definicion}[Autómata de misión]
\label{def:automata}
Para cada carga $\ell$, el modo
$m_\ell\in\{\textsf{REC},\textsf{FORM},\textsf{TRANS},\textsf{DELIV}\}$ evoluciona
según
\[
  \textsf{REC}\ \xrightarrow{\ \text{coalición certificada}\ }\ \textsf{FORM}
  \ \xrightarrow{\ \text{contacto y autoridad}\ }\ \textsf{TRANS}
  \ \xrightarrow{\ \text{demanda extinta}\ }\ \textsf{DELIV},
\]
con una transición de excepción a \textsf{RESC} desde cualquier fase ante fallo,
y retorno a \textsf{TRANS} tras sustitución. Las guardas son, respectivamente:
el margen de \cref{th:robusto} junto con el span positivo de \cref{th:span};
la compatibilidad cinemática de \cref{pr:rigidez} junto con la guarda espectral
de \cref{co:guarda}; y la extinción de \cref{pr:extincion}.
\end{definicion}

\figAutomata

\begin{proposicion}[Los conjuntos admisibles se anidan]
\label{pr:anidados}
Los conjuntos admisibles de las tres fases satisfacen
\[
  \Kcal_{\textsf{REC}}\ \supseteq\ \Kcal_{\textsf{FORM}}\ \supseteq\ \Kcal_{\textsf{TRANS}} .
\]
\end{proposicion}

\begin{proof}
Al pasar a \textsf{FORM} se añaden las restricciones de aproximación al puesto y
las cotas de esfuerzo de \cref{pr:cota}. Al pasar a \textsf{TRANS} se añaden las
igualdades de rigidez de \cref{th:rigidizacion} —que sustituyen desigualdades
internas por igualdades, restringiendo estrictamente— y las barreras de la
envolvente compuesta de \cref{pr:barrera-compuesta}. Añadir restricciones no
amplía el conjunto.
\end{proof}

\Cref{pr:anidados} explica por qué la conmutación es delicada en un solo
sentido: avanzar de fase siempre es admisible si la guarda se cumple, porque el
estado ya pertenece al conjunto mayor; retroceder —disolver una coalición
rígida— sale del conjunto pequeño y requiere la comprobación de holgura de
\cref{pr:rigidizacion-admisible}.

\subsection{Compatibilidad cinemática de la formación}

En modo \textsf{cargo}, el acoplamiento rígido impide deslizamiento lateral del
contacto: el punto material de la carga que coincide con el AMR $i$ debe moverse
a lo largo del rumbo de este.

\begin{proposicion}[Restricción de rigidez]
\label{pr:rigidez}
Con $r_i$ el vector de contacto en el marco de la carga y
$\xi_\ell=(V_x,V_y,\Omega)$ su \emph{twist}, la ausencia de deslizamiento
lateral en todos los contactos equivale a $A_{\mathrm{kin}}\xi_\ell=0$ con filas
\[
  [A_{\mathrm{kin}}]_i=\bigl(-\sin\theta_i,\ \cos\theta_i,\ e_i^{\!\top}r_i\bigr).
\]
La dimensión del espacio de maniobras rígidas admisibles es
$3-\operatorname{rank}A_{\mathrm{kin}}$, y una maniobra deseada $\xi^d_\ell$ es
realizable sin deslizamiento si y solo si $A_{\mathrm{kin}}\xi^d_\ell=0$.
\end{proposicion}

\begin{proof}
La velocidad del punto material en $r_i$ es $V_\ell+\Omega\,Jr_i$ con $J$ la
rotación de $\pi/2$. La condición de no deslizamiento lateral es
$e_i^{\perp\top}(V_\ell+\Omega Jr_i)=0$. Desarrollando,
$e_i^{\perp\top}V_\ell=-\sin\theta_i V_x+\cos\theta_i V_y$ y
$e_i^{\perp\top}Jr_i=e_i^{\!\top}r_i$, lo que da la fila enunciada.
\end{proof}

\begin{observacion}[De dónde sale la incompatibilidad de orientación]
\label{obs:incompatible}
Si $\operatorname{rank}A_{\mathrm{kin}}=3$ el único \emph{twist} sin
deslizamiento es el nulo: la coalición está bloqueada aunque cada AMR tenga
autoridad individual. Y si $\sigma_{\min}(A_{\mathrm{kin}})$ es pequeño, el
\emph{twist} admisible existe pero exige esfuerzos que crecen como su inverso.
Esta es la misma estructura de \cref{th:degradacion}: la incompatibilidad de
orientación no es un caso aparte, es autoridad degradada por geometría de
formación en lugar de por geometría de contacto.
\end{observacion}

\begin{observacion}[Alcance de la restricción: el contacto no está en el eje]
\label{obs:rigidez-alcance}
\Cref{pr:rigidez} deriva $A_{\mathrm{kin}}$ de exigir que el punto de contacto se
mueva a lo largo del rumbo del AMR, lo que sitúa implícitamente el contacto en el
centro del eje de tracción. Físicamente no es así: el pusher o la pinza están
adelantados una distancia $a>0$, y por \cref{pr:fuera-eje} un punto fuera del eje
admite \emph{cualquier} velocidad planar. La igualdad
$A_{\mathrm{kin}}\zeta_\ell=0$ deja entonces de ser una restricción sobre el
\emph{twist} de la carga.

Lo que queda en su lugar es una condición de saturación: para realizar
$\zeta_\ell$, cada AMR necesita
$(u_i,\omega_i)=M(\psi_i,a_i)^{-1}\bigl(V_\ell+\Omega J r_i\bigr)$, y esas
demandas deben caber en sus cotas de \cref{pr:cota}, con la demanda angular
creciendo como $1/a_i$.

La corrección mejora el resultado en vez de debilitarlo. El bloqueo por rango
—$\operatorname{rank}A_{\mathrm{kin}}=3$ y sólo el \emph{twist} nulo admisible—
es un artefacto de $a=0$; con contactos reales el fenómeno no desaparece, pero
cambia de naturaleza: pasa de imposibilidad discreta a saturación continua,
gobernada por la misma ley $1/a$ de \cref{obs:fuera-eje-degradacion}. Una
coalición mal orientada no queda bloqueada, queda lenta, y eso es exactamente lo
que \cref{th:degradacion} predice para todos los demás recursos. La guarda de
\cref{def:automata} debe leerse, en consecuencia, sobre la saturación y no sobre
el rango.
\end{observacion}

\subsection{Frenado admisible: descargar la hipótesis de factibilidad}

\Cref{pr:hocbf} garantiza invariancia \emph{si el programa es factible en todo
estado alcanzado}, y esa hipótesis quedó sin descargar. La construcción
siguiente la convierte en propiedad.

\begin{lema}[Cota de velocidad por frenado]
\label{le:astop}
Sea $a^{\max}_i$ la deceleración máxima del AMR $i$ y $h_{\min}$ la holgura
mínima a un obstáculo o vecino. Si se impone como barrera adicional
\[
  v_i\ \le\ \sqrt{2\,a^{\max}_i\,(h_{\min}-d_{\min})},
\]
entonces la maniobra de frenado a deceleración máxima es siempre admisible, y en
consecuencia el conjunto admisible del programa de movimiento es no vacío desde
todo estado alcanzable.
\end{lema}

\begin{proof}
La distancia de frenado desde $v_i$ es $v_i^{2}/(2a^{\max}_i)$, que bajo la cota
enunciada no excede $h_{\min}-d_{\min}$. Frenar es por tanto una acción que
respeta todas las barreras, luego el conjunto admisible las contiene.
\end{proof}

\begin{corolario}[Invariancia incondicional]
\label{co:invariancia-incondicional}
Bajo \cref{le:astop} y con prioridad lexicográfica de las restricciones de
seguridad en el programa, la conclusión de \cref{pr:hocbf} se sostiene sin la
hipótesis de factibilidad: el conjunto seguro es invariante desde todo estado
inicial admisible, con independencia de lo que decidan las capas superiores.
\end{corolario}

\begin{observacion}[El precio de la garantía]
\label{obs:precio-astop}
La cota de \cref{le:astop} es restrictiva cuando la holgura es escasa: con
$a^{\max}=1{,}2$~m/s$^2$ y $d_{\min}=0{,}5$~m, una holgura de $0{,}6$~m limita a
$0{,}49$~m/s, mientras que $2{,}0$~m permite $1{,}90$~m/s. La velocidad de
operación queda determinada por la holgura disponible, no elegida. Esa pérdida
de caudal entra en $\Delta_{\mathrm{dyn}}$ y compite con la cota de grado de
\cref{pr:compatibilidad}: pasillos estrechos limitan a la vez la velocidad y el
árbol de conectividad.
\end{observacion}

\subsection{Frenar el cuerpo no es frenar los AMR}

\Cref{le:astop} acota la velocidad de cada AMR por su propia deceleración. En
\textsf{TRANS} eso ya no basta: quien debe detenerse es el cuerpo compuesto, cuya
masa incluye la carga y cuya capacidad de frenado no es la suma de los motores
sino lo que el conjunto de \emph{wrench} realizable permite en la dirección
contraria al avance.

\begin{teorema}[Frenado del cuerpo compuesto]
\label{th:astop-cuerpo}
Sea $m_{\mathrm{tot}}=m_\ell+\sum_{i\in\Ccal_\ell}m_i$, sea $\hat v$ la dirección
del movimiento del cuerpo y sea
\[
  F_{\mathrm{fren}}(\hat v)\;=\;h_{\mathcal W(\Ccal_\ell)}(-\hat v)
  \;=\;\max_{w\in\mathcal W(\Ccal_\ell)}\bigl(-\hat v^{\!\top}w^{\mathrm{lin}}\bigr)
\]
la función soporte del conjunto de \emph{wrench} realizable en la dirección de
frenado. Supóngase $F_{\mathrm{fren}}$ constante a lo largo de la maniobra, o
sustitúyase por su mínimo sobre las direcciones barridas. Si se impone la barrera
\begin{equation}
  h_{\mathrm{cuerpo}}
  \;=\;h-d_{\min}-\frac{m_{\mathrm{tot}}\lVert V_\ell\rVert^{2}}
  {2\,F_{\mathrm{fren}}(\hat v)}\ \ge\ 0,
  \label{eq:barrera-cuerpo}
\end{equation}
entonces desde todo estado con $h_{\mathrm{cuerpo}}\ge0$ el frenado a capacidad
máxima es realizable y $h_{\mathrm{cuerpo}}$ se mantiene no negativa: el cuerpo
se detiene dentro de su holgura.
\end{teorema}

\begin{proof}
Factibilidad: por definición de función soporte existe
$w\in\mathcal W(\Ccal_\ell)$ con $-\hat v^{\!\top}w^{\mathrm{lin}}=F_{\mathrm{fren}}$, y
ese \emph{wrench} imprime la deceleración $F_{\mathrm{fren}}/m_{\mathrm{tot}}$.
Invariancia: con ese control $\dot{\lVert V_\ell\rVert}=-F_{\mathrm{fren}}/m_{\mathrm{tot}}$,
luego
\[
  \frac{d}{dt}\Bigl[\frac{m_{\mathrm{tot}}\lVert V_\ell\rVert^{2}}{2F_{\mathrm{fren}}}\Bigr]
  =\frac{m_{\mathrm{tot}}\lVert V_\ell\rVert}{F_{\mathrm{fren}}}\,\dot{\lVert V_\ell\rVert}
  =-\lVert V_\ell\rVert .
\]
Como la holgura no decrece más deprisa que la velocidad de acercamiento,
$\dot h\ge-\lVert V_\ell\rVert$, se obtiene
$\dot h_{\mathrm{cuerpo}}\ge-\lVert V_\ell\rVert+\lVert V_\ell\rVert=0$.
\end{proof}

\begin{observacion}[La capacidad de frenar es un \emph{wrench}]
\label{obs:frenado-soporte}
\Cref{th:astop-cuerpo} no añade una teoría nueva: la reutiliza. El objeto que
mide la capacidad de frenado es la misma función soporte con la que
\cref{th:robusto} mide el margen de realizabilidad, evaluada en otra dirección.
Seguridad del cuerpo y certificado de mercado dejan de ser dos comprobaciones
independientes y pasan a ser dos evaluaciones del mismo conjunto
$\mathcal W(\Ccal_\ell)$. En particular, todo lo que estrecha ese conjunto —fricción
baja, un AMR perdido, contacto mal orientado— reduce a la vez la capacidad de
transportar y la de detenerse.

Una simulación con cuatro contactos aleatorios, $m_{\mathrm{tot}}=400$~kg y
esfuerzos acotados a $150$~N da $F_{\mathrm{fren}}$ entre $106$ y $223$~N, con
velocidades admisibles de $1{,}15$ a $1{,}67$~m/s desde $3$~m de holgura; en las
tres realizaciones el cuerpo se detiene con $h_{\mathrm{cuerpo}}$ tocando cero
sin cruzarlo, que es el comportamiento de una barrera ajustada.
\end{observacion}

\begin{observacion}[Certificar con la autoridad que queda, no con la que había]
\label{obs:certificado-postevento}
El criterio de autorización de una maniobra de frenado ante un obstáculo que
aparece a distancia $d$ es $d\ge d_s+vT_{\mathrm{reac}}+v^2/(2\underline a)$,
y la única decisión es qué $\underline a$ se usa. En mil frenados
longitudinales de una carga de $600$--$1\,400$~kg con cuatro apoyos, velocidad
$0{,}3$--$0{,}7$~m/s, retardo de reacción y una pérdida parcial de autoridad
uniforme entre el $40\%$ y el $100\%$ de la nominal, el certificado con la
autoridad anterior al evento autorizó $809$ maniobras de las que $176$ no eran
seguras —$21{,}8\%$ de falsas aceptaciones, intervalo $19{,}0$--$24{,}8\%$—;
el certificado con la cota inferior válida de la autoridad posterior autorizó
$303$ sin ninguna falsa aceptación, con cota superior de Clopper--Pearson del
$1{,}21\%$, y rechazó $330$ maniobras que en realidad eran seguras
\cite{mrob2026r11}. Tres lecturas. La cota es conservadora por construcción,
que es el mismo precio que \cref{pr:recourse} y \cref{pr:commit}. Rechazar un
certificado no crea una maniobra de rescate: $367$ episodios estaban fuera de
la región de frenado seguro con cualquier certificado, y la barrera de
\cref{th:astop-cuerpo} es la que debe haber evitado llegar ahí. Y el mismo
criterio aplicado por cualquier otro planificador produce las mismas mil
decisiones: la ventaja es usar una cota física válida tras el evento, no una
propiedad del juego.
\end{observacion}


\begin{observacion}[Cuál de las dos cotas manda]
\label{obs:dos-frenados}
Las cotas de \cref{le:astop} y \cref{th:astop-cuerpo} no son comparables a
priori y deben imponerse ambas. En \textsf{REC} y \textsf{FORM} solo aplica la
primera, porque no hay cuerpo. En \textsf{TRANS} suele mandar la segunda, porque
$m_{\mathrm{tot}}$ es mayor que la masa del AMR en un factor grande mientras que
$F_{\mathrm{fren}}$ está limitada por la fricción disponible y no por los
motores. Ese cambio de cota al cruzar la guarda de \textsf{FORM} a
\textsf{TRANS} es una discontinuidad de velocidad admisible que el planificador
debe anticipar: la coalición ha de frenar \emph{antes} de acoplarse, no después.
\end{observacion}

\subsection{La entrega no necesita orden}

\begin{proposicion}[Extinción del precio y disolución]
\label{pr:extincion}
Sea la demanda de \emph{wrench} generada por el seguimiento,
$\Wdem_\ell=M_\ell\dot\xi^d_\ell+D_\ell\xi^d_\ell$. Si al aproximarse al destino
$\xi^d_\ell\to0$, entonces $\Wdem_\ell\to0$; el precio del ajuste
\eqref{eq:tatonnement} converge a cero, y con él el pago por participar. En
consecuencia la coalición se disuelve por la misma dinámica de revisión que la
formó, sin orden explícita de liberación.
\end{proposicion}

\begin{proof}
Con $\Wdem_\ell\to0$ el punto de reposo del ajuste con fuga $d_p>0$ es
$p_\ell=0$. Por \cref{th:potencial}, el pago de pertenecer a la coalición
incluye el término proporcional al precio; anulado este, la desviación a
\textsf{libre} deja de estar penalizada y el potencial deja de premiar la
permanencia.
\end{proof}

\begin{observacion}[Comprobación numérica]
\label{obs:extincion-num}
Con $\xi^d$ decreciendo con la distancia al destino, el ajuste de precios da
$\lVert p\rVert=0{,}999$ a $5$~m, $0{,}499$ a $0{,}5$~m, $0{,}050$ a $0{,}05$~m y
exactamente cero en el destino. La disolución es gradual y proporcional, no un
evento discreto: los esfuerzos de contacto se relajan antes de que el cuerpo
compuesto deje de existir.
\end{observacion}

\subsection{Ejecutar mientras se sigue negociando}

Hay una carencia en todo lo anterior que los ensayos de \cref{sec:ensayos}
exhiben sin ambigüedad. El autómata de \cref{def:automata} encadena fases y
\cref{th:multiclase} tarifa rutas, pero ninguno describe lo que el banco hace en
realidad: recalcular propuestas de trayectoria cada $0{,}4$~s y ejecutar sin
detenerse mientras el optimizador sigue trabajando. Esa capa merece su propio
enunciado, porque tiene garantías propias.

\begin{definicion}[Testigo vigente]
\label{def:testigo}
Un \emph{testigo} es un plan completo desde el conjunto de creencia actual hasta
el objetivo, que incluye un tubo válido ante las perturbaciones declaradas, las
guardas y reinicializaciones seguras, una ocupación temporal coherente y una
política de ejecución local, con la salida de cada tramo contenida en la entrada
del siguiente. Su coste restante admite una envolvente $\Vcal$ que decrece al
menos $\ell_0\Delta t>0$ por unidad de tiempo físico fuera del objetivo.
\end{definicion}

\begin{teorema}[Seguridad y progreso con testigo vigente]
\label{th:testigo}
Supóngase que existe un testigo completo inicial de coste $\Vcal_0$, que la
ejecución no se detiene durante el cálculo —sigue el testigo vigente—, que toda
sustitución se confirma con la misma versión e instante de activación dentro de
un tubo de entrada válido, y que no hay Zenón. Admítanse únicamente las
sustituciones cuyo coste restante, incluida la transición física, no exceda al
del testigo vigente, y cuéntense como \emph{estrictas} las que además descienden
al menos $\delta>0$ \cite{mayorga2026m1}. Entonces la ejecución es segura, alcanza el objetivo en
tiempo no mayor que $\Vcal_0/\ell_0$, su coste acumulado no supera $\Vcal_0$, y
el número de revisiones estrictas no excede $\lfloor\Vcal_0/\delta\rfloor$.

En particular, \emph{no se exige que el optimizador no lineal converja antes de
cada mando}.
\end{teorema}

\begin{observacion}[Comprobación de las tres cotas]
\label{obs:testigo-num}
Sobre ocho realizaciones con $\Vcal_0=100$, $\ell_0=2$ y $\delta=5$ —cotas
teóricas de $50$~s, coste $100$ y $20$ revisiones estrictas— los tiempos medidos
van de $10{,}0$ a $20{,}0$~s, los costes de $26{,}1$ a $52{,}1$, y las revisiones
estrictas de $4$ a $8$. Las tres cotas se respetan en todas las realizaciones, y
las propuestas que empeoran el coste restante se rechazan sin detener la
ejecución.
\end{observacion}

\begin{observacion}[Qué cambia esto respecto del resto del documento]
\label{obs:testigo-alcance}
Tres cosas. Primera, da la cota de tiempo que \cref{obs:complejidad-propia}
declaraba ausente, aunque solo para la capa de trayectoria y en función de
$\Vcal_0$ y $\ell_0$, que dependen de la instancia. Segunda, explica por qué el
ensayo P01 puede sostener $131$ épocas de negociación sin detenerse ni una vez
dentro del pasillo: la negociación no es un preludio de la ejecución sino
concurrente con ella. Tercera, invierte la relación habitual con la
optimalidad: no se busca el óptimo y luego se ejecuta, sino que se ejecuta algo
admisible y se acepta toda mejora que llegue, con el número de mejoras acotado
por el presupuesto inicial.

La hipótesis crítica es la existencia del testigo inicial, que en este trabajo
proporciona el certificado de \cref{th:robusto} junto con la factibilidad de
frenado de \cref{le:astop}. Y hay una salvedad energética que conviene no
esconder en un coste terminal: una batería finita con consumo auxiliar positivo
no permite sostener indefinidamente la condición terminal, de modo que el
mantenimiento seguro del objetivo debe declararse con horizonte o con recarga,
que es lo que hace \cref{pr:caudal}.
\end{observacion}

\subsection{La versión esperada: progreso con residual de Bellman}

\Cref{th:testigo} es determinista y exige que cada sustitución no aumente el
coste restante. Su contrapartida en esperanza admite un residual por unidad de
duración, y es la que corresponde a un catálogo con transiciones inciertas
\cite{mayorga2026m1}.

\begin{teorema}[Progreso esperado con residual por unidad de duración]
\label{th:bellman-residual}
Sea $W\ge0$ una función de valor finita y fija durante la ejecución, con
$W(g)=0$ en el objetivo. En cada estado no terminal se ejecuta una continuación
$b$ de duración $\tau\ge\tau_{\min}>0$ y coste esperado
$c(z,b)\ge\ell_0\,\mathbb E[\tau\mid z,b]$, $\ell_0>0$. Si toda continuación
seleccionada satisface
\[
  c(z,b)+\mathbb E\bigl[W(z')\mid z,b\bigr]-W(z)\ \le\ \varrho\,\mathbb E[\tau\mid z,b],
  \qquad 0\le\varrho<\ell_0,
\]
y el sistema puede continuar hasta el objetivo, entonces
\[
  \mathbb E[T_g]\ \le\ \frac{W(z_0)}{\ell_0-\varrho},
  \qquad
  \mathbb E[J]\ \le\ \frac{\ell_0}{\ell_0-\varrho}\,W(z_0),
\]
y sin eventos de Zenón el objetivo se alcanza casi seguramente.
\end{teorema}

\begin{proof}
Sumando la desigualdad a lo largo de las decisiones hasta $n\wedge N_g$ y usando
esperanza iterada, $\mathbb E\sum c_k\le W(z_0)-\mathbb E\,W(z_{n\wedge N_g})
+\varrho\,\mathbb E\sum\tau_k$. Como $W\ge0$ y $c_k\ge\ell_0\mathbb E[\tau_k]$,
resulta $(\ell_0-\varrho)\mathbb E\sum\tau_k\le W(z_0)$; el límite por
convergencia monótona da ambas cotas. Una probabilidad positiva de infinitas
decisiones sin entrega haría infinito el tiempo esperado por $\tau\ge\tau_{\min}$,
contradicción.
\end{proof}

\begin{observacion}[Comprobación y alcance]
\label{obs:bellman-num}
Sobre una cadena de seis estados con dos continuaciones por estado —lenta y
rápida, de duraciones $2{,}0$ y $0{,}8$— y veinte mil realizaciones, con
$W(z_0)=12$ y $\ell_0=1$: $\mathbb E[T_g]=8{,}75$ frente a la cota $12$, y
$\mathbb E[J]=10{,}06$ frente a $12$. No es una cota determinista del peor
tiempo, y la seguridad física \emph{no} se hereda: solo si cada continuación y
sus transiciones son seguras por separado, que es lo que \cref{def:testigo}
exige. Su utilidad es permitir error numérico controlado sin resolver
exactamente cada juego: $\varrho$ es el residuo de Bellman que se tolera por
segundo de maniobra.
\end{observacion}

\begin{proposicion}[Presupuesto de error de una evaluación local]
\label{pr:presupuesto-error}
Sean $R=c+PW-W(z)$ el residual verdadero y $\hat R=\hat c+\hat P\hat W-\hat W(z)$
el calculado con información local. Si $|\hat c-c|\le\epsilon_c$,
$\lVert\hat W-W\rVert_\infty\le\epsilon_V$, $\lVert\hat P-P\rVert_1\le\epsilon_P$ y
$\lVert W\rVert_\infty\le M_W$, entonces
$|\hat R-R|\le\epsilon_c+2\epsilon_V+M_W\epsilon_P$.
\end{proposicion}

\begin{observacion}[Cuándo pedir más información y cuándo seguir]
\label{obs:presupuesto-uso}
La regla $\hat R+\epsilon_R\le\varrho\,\underline\tau$, con
$\underline\tau\le\mathbb E[\tau]$, es suficiente para la hipótesis de
\cref{th:bellman-residual}. Eso convierte al presupuesto en una decisión
operativa: si con la información disponible el residual certificado cabe en
$\varrho\underline\tau$, se ejecuta; si no, se solicita más información o se
sigue el respaldo vigente. Es el mismo criterio de \cref{th:intervalo} llevado
del valor de una rama al progreso de la misión, y enlaza con \cref{th:saltos}:
cada salto adicional reduce $\epsilon_V$ y $\epsilon_P$ a un ritmo que fija el
condicionamiento. La proposición no proporciona por sí sola las cotas de
cartografía, contacto o probabilidad —$\epsilon_V,\epsilon_P$ deben
identificarse y validarse—, y así se declara.
\end{observacion}

\begin{observacion}[Qué valor debe comparar el reclutamiento]
\label{obs:equipo-barato}
Un ejemplo sin solver fija la moraleja. Dos continuaciones cuestan $1$ y $4$
unidades; la primera conduce a un estado cuya mejor entrega restante cuesta
$100$, la segunda a uno que cuesta $5$. El criterio inmediato elige la primera y
paga $101$; el de Bellman elige la segunda y paga $9$. El equipo barato al
principio puede ser caro hasta el final, y el reclutamiento debe comparar
coste hasta la entrega, no coste de la maniobra. Dos salvedades más: con
criterio descontado en tiempo físico hay que usar el factor de duración
$\mathbb E[e^{-\beta\tau}W(z')]$ y no un descuento idéntico para maniobras de
distinta duración; y una abstracción por pose sola no suele ser markoviana
cuando contactos, desgaste, información o reservas tienen memoria, de modo que
hay que ampliar el estado, usar una creencia, o acotar el error de abstracción.
\end{observacion}

\subsection{Qué garantiza el encadenamiento y qué no}

\begin{teorema}[Propiedades de la misión encadenada]
\label{th:mision}
Supóngase: admisión certificada en cada reclutamiento; conectividad factorizada
(\cref{pr:factorizada}); guarda de compatibilidad (\cref{pr:rigidez}) y
espectral (\cref{co:guarda}) en la adquisición; cota de frenado
(\cref{le:astop}); y ganancia pequeña entre capas (\cref{th:smallgain}).
Entonces:
\begin{enumerate}[label=(\roman*),leftmargin=2.2em,topsep=0.2em,itemsep=0.15em]
  \item el estado permanece en el conjunto seguro en todo instante y en toda
  fase, incondicionalmente (\cref{co:invariancia-incondicional});
  \item cada transición de fase está gobernada por una guarda verificable
  localmente, y la frecuencia de conmutación está acotada por
  \cref{th:dwell};
  \item la degradación de autoridad se detecta antes de sufrirse, porque la
  guarda espectral se evalúa sobre el estado predicho (\cref{co:guarda});
  \item la entrega ocurre por extinción del precio, sin coordinación adicional
  (\cref{pr:extincion}).
\end{enumerate}
\end{teorema}

\begin{proof}
(i) es \cref{co:invariancia-incondicional} aplicado en cada fase, y las
restricciones se anidan por \cref{pr:anidados}, luego la invariancia de la fase
más restrictiva implica la de las anteriores. (ii) reúne las guardas de
\cref{def:automata} con el tiempo de permanencia. (iii) y (iv) son
\cref{co:guarda} y \cref{pr:extincion}.
\end{proof}

\begin{observacion}[Lo que este teorema no afirma]
\label{obs:mision-limites}
No se afirma optimalidad de la misión completa. En particular, \emph{no} se
sostiene que el coste total se descomponga en la suma de los óptimos por fase
más términos de interfaz acotados: cada fase resuelve su problema tomando como
dado el resultado de la anterior, y esa secuencialidad es precisamente una de
las fuentes de $\Delta_{\mathrm{dyn}}$ en \cref{th:descomposicion}. Tampoco se
afirman cotas de tiempo por fase: acotar $T_{\textsf{REC}}$ exige la tasa de
mezcla efectiva del canal, y acotar $T_{\textsf{TRANS}}$ exige el número de
conflictos, que depende de la instancia. Lo que el encadenamiento garantiza es
seguridad, verificabilidad de las guardas y terminación de la disolución; no
rendimiento.
\end{observacion}
